%% kapa: exam=1 round=28 year=2017 period=2 form=A subject=accounting subject_ko=회계학 seq=18 number=58
%% source: Q-Net 감정평가사 기출문제 (cst00309) artlSeq=5199603 file=extracted/1008/5199603/2017년 제28회 감정평가사 1차시험 문제지/2교시 A형.pdf
%% conversion: kordoc 4.24.51 -> tools/gen1.py (auto); review=proofread-sonnet; flags=none
\begin{question}{58}
재무제표 표시에 관한 설명으로 옳은 것은?
\begin{choices}
  \item 비용을 기능별로 분류하는 것이 성격별 분리보다 더욱 목적적합한 정보를 제공하므로, 비용은 기능별로 분류한다.
  \item 재무상태표에 표시되는 자산과 부채는 반드시 유동자산과 비유동자산, 유동부채와 비유동부채로 구분하여 표시하여야 한다.
  \item 영업이익에 포함되지 않은 항목 중 기업의 영업성과를 반영하는 그 밖의 수익 항목이 있다면 조정영업이익으로 포괄손익계산서 본문에 표시하여야 한다.
  \item 재무제표에는 중요하지 않아 구분하여 표시하지 않은 항목이라도 주석에서는 구분 표시해야 할 만큼 충분히 중요할 수 있다.
  \item 부적절한 회계정책은 이에 대하여 공시나 주석 또는 보충자료를 통해 설명할 수 있다면 정당화 될 수 있다.
\end{choices}
\end{question}
\answer{4}
